\documentclass[a4paper,11pt]{article}
\usepackage[utf8]{inputenc}
\usepackage[usenames,dvipsnames]{color}
\usepackage{graphicx}
\usepackage[justification=centering,labelfont=bf]{caption}
\usepackage{listings}
\begin{document}
\begin{titlepage}
\begin{center}
\textsc{\Large Parallelism}
\\
\texttt{1202}
\\[1.5cm]
\rule{\linewidth}{0.5mm}
\\[0.4cm]
{\huge
\bfseries
Lab 1: Embarrassingly parallelism with OpenMP: Mandelbrot set
\\[0.4cm]
}
\rule{\linewidth}{0.5mm}
\\[2.5cm]
\begin{minipage}{0.4\textwidth}
\begin{flushleft}
\large
Héctor Ramón Jiménez
\end{flushleft}
\end{minipage}
\begin{minipage}{0.4\textwidth}
\begin{flushright}
\large
Alvaro Espuña Buxo
\end{flushright}
\end{minipage}
\vfill
{\large
\today
}
\\
{\large
\texttt{Facultat d'Informàtica de Barcelona}
}
\end{center}
\end{titlepage}
\section{Reading activity}
In this first part of the report we first briefly describe the basic formulation for the \textbf{Mandelbrot set}
    (reference here) and then its implementation in the sequential code available (\texttt{mandel-serial.c}).
\subsection{Description}
The \textbf{Mandelbrot set} is a particular set of points in the complex domain named after the mathematician
    \textbf{Benoit Mandelbrot}, who studied it and popularized it. The \textbf{Mandelbrot set}'s boundary is an easily
    recognizable two-dimensional fractal shape (see figure \ref{mandelbrot:boundary}).
\begin{figure}[h!]
\includegraphics[width=1.0\textwidth]{figures/mandelbrot_boundary.jpg}
\caption{The Mandelbrot set boundary}
\label{mandelbrot:boundary}
\end{figure}
\\
Mandelbrot set images are made by sampling complex numbers and determining for each whether the result
    \textbf{tends towards infinity} when a particular mathematical operation is \textbf{iterated} on it. The real
    and imaginary parts of each number are treated as image coordinates. The pixels are colored according to how rapidly
    the sequence diverges.
\\
More precisely, the \textbf{Mandelbrot set} is the set of values of $c$ in the complex plane for which the orbit of
    $0$ under iteration of the complex quadratic polynomial $z_{n+1} = z^2_n + c$ remains bounded.
\subsection{Algorithm}
The simplest algorithm for drawing a picture of the Mandelbrot set is the following. We discretize the
    complex space in a set of points and each point corresponds to a pixel in a two dimensional plot. To color
    any such pixel, let c (represented by the complex variable ...... [Complete it] at the mandel-serial.c
    code) be the midpoint of that pixel. ... [ Here you should continue the description of the sequential
    algorithm clearly naming the variables that are used in mandel-serial.c and the criteria used to decide
    the color for each point.]
\\
Here you should insert the identification reference and the description of that reference: title,
    authors, year of publication, website, etc.
\end{document}
